\section{\ENG{Metadata} \ENG{Standards} \ENG{Foundation} (1995-2002)}
\label{sec:metadata_foundation}

Παράλληλα με την ανάπτυξη των προτύπων επικοινωνίας και δομής περιεχομένου (όπως το \ENG{AICC}), η δεκαετία του 1990 σηματοδότησε την έναρξη συστηματικών προσπαθειών για την τυποποίηση των μεταδεδομένων (\textit{\ENG{metadata}}) εκπαιδευτικών πόρων. Τα \ENG{metadata} αποτελούν «δεδομένα για δεδομένα» και περιγράφουν χαρακτηριστικά του εκπαιδευτικού περιεχομένου όπως ο τίτλος, ο συγγραφέας, η γλώσσα, το επίπεδο δυσκολίας, οι μαθησιακοί στόχοι και τεχνικά χαρακτηριστικά~\cite{dcmi2020dcmi_terms,ieee2002lom_standard,dublincore2024metadata}. Η ύπαρξη τυποποιημένων μεταδεδομένων είναι κρίσιμη για την ανευρεσιμότητα (\textit{\ENG{discoverability}}), την επαναχρησιμοποιησιμότητα (\textit{\ENG{reusability}}) και την διαλειτουργικότητα (\textit{\ENG{interoperability}}) εκπαιδευτικών πόρων.

Σε πρακτικό επίπεδο, τα \ENG{metadata standards} υποστηρίζουν τη διασύνδεση διαφορετικών συστημάτων και την ανταλλαγή εκπαιδευτικών πόρων μέσω \ENG{repositories}, διευκολύνοντας την αναζήτηση και την επαναχρησιμοποίηση σε ετερογενή περιβάλλοντα~\cite{inokufu2024metadata,researchgate2025metadata}.

\subsection{\ENG{Dublin} \ENG{Core} \ENG{Metadata} \ENG{Initiative} (1995)}
\label{subsec:dublin_core}

Το \ENG{Dublin} \ENG{Core} \ENG{Metadata} \ENG{Element} \ENG{Set} αναπτύχθηκε το 1995 κατά τη διάρκεια ενός εργαστηρίου στο \ENG{Dublin} του \ENG{Ohio}, με στόχο τη δημιουργία ενός απλού και ευρέως εφαρμόσιμου συνόλου μεταδεδομένων για την περιγραφή ψηφιακών πόρων~\cite{dcmi2020dcmi_terms,dcpapers2024educational}. Το \ENG{Dublin} \ENG{Core} \ENG{Metadata} \ENG{Initiative} (\ENG{DCMI}) δημιούργησε ένα σύνολο 15 βασικών στοιχείων μεταδεδομένων που μπορούν να εφαρμοστούν σε οποιονδήποτε τύπο ψηφιακού πόρου, συμπεριλαμβανομένου του εκπαιδευτικού περιεχομένου.

Τα 15 βασικά στοιχεία του \ENG{Dublin} \ENG{Core} είναι:

\begin{table}[H]
\centering
\caption{Βασικά Στοιχεία \ENG{Dublin} \ENG{Core} \ENG{Metadata}}
\label{tab:dublin_core_elements}
\small
\setlength{\tabcolsep}{5pt}
\renewcommand{\arraystretch}{1.25}
\begin{chaptertablebox}
\begin{tabular}{|l|p{8cm}|}
\hline
\textbf{Στοιχείο} & \textbf{Περιγραφή} \\
\hline
\hline
\ENG{Title} & Το όνομα που δίνεται στον πόρο \\
\ENG{Creator} & Η οντότητα που είναι κυρίως υπεύθυνη για τη δημιουργία του πόρου \\
\ENG{Subject} & Το θέμα του περιεχομένου του πόρου \\
\ENG{Description} & Μια περιγραφή του περιεχομένου του πόρου \\
\ENG{Publisher} & Η οντότητα που είναι υπεύθυνη για τη διάθεση του πόρου \\
\ENG{Contributor} & Οντότητα που είναι υπεύθυνη για συνεισφορές στο περιεχόμενο \\
\ENG{Date} & Ημερομηνία που σχετίζεται με το κύκλο ζωής του πόρου \\
\ENG{Type} & Η φύση ή το είδος του περιεχομένου \\
\ENG{Format} & Το φυσικό ή ψηφιακό μορφότυπο του πόρου \\
\ENG{Identifier} & Μια μοναδική αναφορά στον πόρο (π.χ. \ENG{URI}, \ENG{ISBN}) \\
\ENG{Source} & Πόρος από τον οποίο προέρχεται ο παρών πόρος \\
\ENG{Language} & Η γλώσσα του περιεχομένου του πόρου \\
\ENG{Relation} & Αναφορά σε σχετικό πόρο \\
\ENG{Coverage} & Το χωρικό ή χρονικό πεδίο εφαρμογής του πόρου \\
\ENG{Rights} & Πληροφορίες για δικαιώματα που σχετίζονται με τον πόρο \\
\hline
\end{tabular}
\end{chaptertablebox}
\end{table}



\subsection{\ENG{IEEE} \ENG{Learning} \ENG{Object} \ENG{Metadata} (\ENG{LOM}) (1997-2002)}
\label{subsec:ieee_lom}

Για να αντιμετωπιστούν οι περιορισμοί του \ENG{Dublin} \ENG{Core} στον τομέα του \ENG{E-Learning}, το \ENG{IEEE} \ENG{Learning} \ENG{Technology} \ENG{Standards} \ENG{Committee} (\ENG{LTSC}) ανέπτυξε το \ENG{IEEE} \ENG{Learning} \ENG{Object} \ENG{Metadata} (\ENG{LOM}), το οποίο εγκρίθηκε ως \ENG{standard} το 2002 με κωδικό \ENG{IEEE} 1484.12.1-2002~\cite{ieee2002lom_standard,ieee2002lom,wiki_lom}.
Το 2020 δημοσιεύτηκε αναθεωρημένη έκδοση ως \ENG{IEEE} 1484.12.1-2020, διατηρώντας τη βασική δομή αλλά επικαιροποιώντας το πλαίσιο περιγραφής~\cite{ieee2020lom}.

Το \ENG{IEEE} \ENG{LOM} είναι ένα εκτεταμένο σχήμα \ENG{metadata} που περιλαμβάνει 76 στοιχεία οργανωμένα σε 9 κατηγορίες:

\begin{table}[H]
\centering
\caption{Οι 9 κατηγορίες του \ENG{IEEE} \ENG{LOM}}
\label{tab:ieee_lom_categories}
\small
\setlength{\tabcolsep}{5pt}
\renewcommand{\arraystretch}{1.25}
\begin{chaptertablebox}
\begin{tabular}{|p{3.2cm}|p{3.8cm}|p{7.2cm}|}
\hline
\textbf{Κατηγορία} & \textbf{Περιγραφή} & \textbf{Ενδεικτικά στοιχεία} \\
\hline
\hline
\textbf{\ENG{General} (Γενικά)} & Γενικές πληροφορίες για τον μαθησιακό πόρο & \ENG{Identifier}, \ENG{Title}, \ENG{Language}, \ENG{Description}, \ENG{Keyword}, \ENG{Coverage}, \ENG{Structure}, \ENG{Aggregation} \ENG{Level} \\
\hline
\textbf{\ENG{Lifecycle} (Κύκλος Ζωής)} & Ιστορικό και κατάσταση του πόρου & \ENG{Version}, \ENG{Status} (\ENG{draft}, \ENG{final}, \ENG{revised}), \ENG{Contribute} (\ENG{authors}, \ENG{publishers}, \ENG{dates}) \\
\hline
\textbf{\ENG{Meta-Metadata}} & Πληροφορίες για τα ίδια τα \ENG{metadata} & \ENG{Identifier}, \ENG{Contribute}, \ENG{Metadata} \ENG{Schema}, \ENG{Language} \\
\hline
\textbf{\ENG{Technical} (Τεχνικά)} & Τεχνικές απαιτήσεις και χαρακτηριστικά & \ENG{Format} (\ENG{MIME} \ENG{type}), \ENG{Size}, \ENG{Location} (\ENG{URL}), \ENG{Requirements} (\ENG{OS}, \ENG{Browser}), \ENG{Installation} \ENG{Remarks}, \ENG{Duration} \\
\hline
\textbf{\ENG{Educational} (Εκπαιδευτικά)} & Παιδαγωγικά χαρακτηριστικά & \ENG{Interactivity} \ENG{Type} (\ENG{active}, \ENG{expositive}, \ENG{mixed})\newline
\ENG{Learning} \ENG{Resource} \ENG{Type} (\ENG{exercise}, \ENG{simulation}, \ENG{questionnaire}, \ENG{diagram}, \ENG{figure}, \ENG{graph}, \ENG{index}, \ENG{slide}, \ENG{table}, \ENG{narrative} \ENG{text}, \ENG{exam}, \ENG{experiment}, \ENG{problem} \ENG{statement}, \ENG{self} \ENG{assessment}, \ENG{lecture})\newline
\ENG{Interactivity} \ENG{Level} (\ENG{very} \ENG{low}, \ENG{low}, \ENG{medium}, \ENG{high}, \ENG{very} \ENG{high})\newline
\ENG{Semantic} \ENG{Density} (\ENG{very} \ENG{low}, \ENG{low}, \ENG{medium}, \ENG{high}, \ENG{very} \ENG{high})\newline
\ENG{Intended} \ENG{End} \ENG{User} \ENG{Role} (\ENG{teacher}, \ENG{author}, \ENG{learner}, \ENG{manager})\newline
\ENG{Context} (\ENG{school}, \ENG{higher} \ENG{education}, \ENG{training}, \ENG{other})\newline
\ENG{Typical} \ENG{Age} \ENG{Range}; \ENG{Difficulty} (\ENG{very} \ENG{easy}, \ENG{easy}, \ENG{medium}, \ENG{difficult}, \ENG{very} \ENG{difficult})\newline
\ENG{Typical} \ENG{Learning} \ENG{Time}; \ENG{Description}, \ENG{Language} \\
\hline
\textbf{\ENG{Rights} (Δικαιώματα)} & Πνευματική ιδιοκτησία και όροι χρήσης & \ENG{Cost} (\ENG{yes}, \ENG{no}), \ENG{Copyright} \ENG{and} \ENG{Other} \ENG{Restrictions}, \ENG{Description} \\
\hline
\textbf{\ENG{Relation} (Σχέσεις)} & Σχέσεις με άλλους μαθησιακούς πόρους & \ENG{Kind} (\ENG{ispartof}, \ENG{haspart}, \ENG{isversionof}, \ENG{hasversion}, \ENG{isformatof}, \ENG{hasformat}, \ENG{references}, \ENG{isreferencedby}, \ENG{isbasedon}, \ENG{isbasisfor}, \ENG{requires}, \ENG{isrequiredby})\newline
\ENG{Resource} (\ENG{Identifier}, \ENG{Description}) \\
\hline
\textbf{\ENG{Annotation} (Σχολιασμοί)} & Σχόλια για εκπαιδευτική χρήση & \ENG{Entity} (\ENG{person} \ENG{or} \ENG{organization}), \ENG{Date}, \ENG{Description} \\
\hline
\textbf{\ENG{Classification} (Ταξινόμηση)} & Ταξινόμηση σε συστήματα κατηγοριοποίησης & \ENG{Purpose} (\ENG{discipline}, \ENG{idea}, \ENG{prerequisite}, \ENG{educational} \ENG{objective}, \ENG{accessibility} \ENG{restrictions}, \ENG{educational} \ENG{level}, \ENG{skill} \ENG{level}, \ENG{security} \ENG{level}, \ENG{competency})\newline
\ENG{Taxon} \ENG{Path}, \ENG{Description}, \ENG{Keyword} \\
\hline
\end{tabular}
\end{chaptertablebox}
\end{table}


\subsubsection{Παράδειγμα \ENG{IEEE} \ENG{LOM} σε \ENG{XML}}

{\selectlanguage{english}\fontencoding{T1}\selectfont

\begin{verbatim}
  <lom xmlns="http://ltsc.ieee.org/xsd/LOM">
  <general>
    <title><string language="en">Introduction to SCORM</string></title>
    <language>en</language>
    <description>
      <string language="en">Basic overview of SCORM standards</string>
    </description>
    <keyword><string language="en">SCORM</string></keyword>
    <keyword><string language="en">E-Learning</string></keyword>
  </general>
  <technical>
    <format>text/html</format>
    <size>2048000</size>
    <location>http://example.com/scorm_intro.zip</location>
    <duration><duration>PT45M</duration></duration>
  </technical>
  <educational>
    <learningResourceType>
      <value>lecture</value>
    </learningResourceType>
    <intendedEndUserRole><value>learner</value></intendedEndUserRole>
    <context><value>higher education</value></context>
    <typicalLearningTime><duration>PT1H</duration></typicalLearningTime>
    <difficulty><value>medium</value></difficulty>
  </educational>
  <rights>
    <cost><value>no</value></cost>
    <copyrightAndOtherRestrictions><value>yes</value></copyrightAndOtherRestrictions>
  </rights>
</lom>
\end{verbatim}
}


Το \ENG{LOM} αξιοποιήθηκε επίσης ως βάση για το \ENG{SCORM} \ENG{metadata profile}, ώστε η περιγραφή των μαθησιακών αντικειμένων να παραμένει συνεπής μεταξύ διαφορετικών \ENG{LMS}~\cite{osu2024metadata}.

\subsection{Υιοθέτηση και Εξελίξεις του \ENG{LOM}}
\label{subsec:lom_adoption}

Το \ENG{IEEE} \ENG{Learning} \ENG{Object} \ENG{Metadata} (\ENG{LOM}) υιοθετήθηκε ευρέως και αποτέλεσε τη βάση για την ανάπτυξη πολλών εθνικών και διεθνών σχημάτων μεταδεδομένων (\ENG{metadata schemas}). Η ευελιξία του μοντέλου επέτρεψε την προσαρμογή του σε διαφορετικά εκπαιδευτικά συστήματα και τεχνολογικά περιβάλλοντα, οδηγώντας στη δημιουργία εξειδικευμένων προφίλ και επεκτάσεων.

Ένα από τα σημαντικότερα παραδείγματα αποτελεί το \ENG{IMS} \ENG{Learning} \ENG{Resource} \ENG{Metadata} (\ENG{LRM}), το οποίο αναπτύχθηκε από τον οργανισμό \ENG{IMS} \ENG{Global}. Το σχήμα αυτό βασίστηκε στο πρότυπο \ENG{IEEE} \ENG{LOM}, αλλά εισήγαγε πρόσθετες επεκτάσεις ώστε να καλύπτει ειδικότερες ανάγκες διαχείρισης και ανταλλαγής εκπαιδευτικών πόρων.

Παράλληλα, στον Καναδά αναπτύχθηκε το \ENG{CanCore}, ένα προφίλ του \ENG{LOM} που παρείχε σαφείς οδηγίες χρήσης και πρακτικές κατευθύνσεις για την εφαρμογή των μεταδεδομένων στο εκπαιδευτικό σύστημα της χώρας. Το \ENG{CanCore} επικεντρώθηκε στη βελτίωση της συνέπειας και της διαλειτουργικότητας μεταξύ διαφορετικών εκπαιδευτικών αποθετηρίων.

Στο Ηνωμένο Βασίλειο δημιουργήθηκε το \ENG{UK} \ENG{LOM} \ENG{Core}, ένα υποσύνολο του \ENG{LOM} που προσαρμόστηκε στις ανάγκες της υποχρεωτικής και της τριτοβάθμιας εκπαίδευσης. Το προφίλ αυτό επικεντρώθηκε στην απλοποίηση του αρχικού μοντέλου και στη διευκόλυνση της υιοθέτησής του από εκπαιδευτικούς οργανισμούς.

Αντίστοιχα, στη Σιγκαπούρη αναπτύχθηκε το \ENG{SingCore}, μια προσαρμογή του \ENG{LOM} που ενσωμάτωσε χαρακτηριστικά προσαρμοσμένα στις ανάγκες του ασιατικού εκπαιδευτικού περιβάλλοντος. Η ύπαρξη αυτών των διαφορετικών προφίλ καταδεικνύει την ευελιξία του \ENG{LOM} και τον ρόλο του ως θεμελιώδους προτύπου για την περιγραφή και την ανταλλαγή μαθησιακών αντικειμένων σε διεθνές επίπεδο.
\subsection{Κριτική και Προκλήσεις του \ENG{IEEE} \ENG{LOM}}
\label{subsec:lom_challenges}

Παρά την ευρεία υιοθέτηση, το \ENG{IEEE} \ENG{LOM} αντιμετώπισε σημαντικές προκλήσεις:

\begin{table}[H]
\centering
\caption{Κύριες προκλήσεις του \ENG{IEEE} \ENG{LOM}}
\label{tab:ieee_lom_challenges}
\small
\setlength{\tabcolsep}{5pt}
\renewcommand{\arraystretch}{1.25}
\begin{chaptertablebox}
\begin{tabular}{|p{3.5cm}|p{9.5cm}|}
\hline
\textbf{Κατηγορία} & \textbf{Προβλήματα} \\
\hline
\hline
\textbf{Πολυπλοκότητα} & 76 στοιχεία - υπερβολική πολυπλοκότητα για \ENG{manual} \ENG{creation}.\newline
Δυσκολία κατανόησης και συμπλήρωσης όλων των πεδίων.\newline
Υψηλό κόστος παραγωγής \ENG{metadata}. \\
\hline
\textbf{Αμφισημία} & Πολλά πεδία έχουν ασαφείς ορισμούς.\newline
Διαφορετικές ερμηνείες από διαφορετικούς \ENG{implementers}.\newline
Έλλειψη συνέπειας μεταξύ \ENG{repositories}. \\
\hline
\textbf{Χαμηλή \ENG{Adoption} \ENG{Rate}} & Οι περισσότεροι \ENG{content} \ENG{creators} δεν συμπλήρωναν πλήρη \ENG{LOM} \ENG{metadata}.\newline
Τα περισσότερα αποθετήρια χρησιμοποιούσαν μόνο μικρό υποσύνολο πεδίων.\newline
\ENG{Automated} \ENG{metadata} \ENG{generation} ήταν περιορισμένη. \\
\hline
\textbf{Ελλιπής Υποστήριξη Εργαλείων} & Περιορισμένα \ENG{authoring} \ENG{tools} με ενσωματωμένη υποστήριξη \ENG{LOM}.\newline
Δύσχρηστα \ENG{metadata} \ENG{editors}.\newline
Περιορισμένη ενσωμάτωση με \ENG{content} \ENG{development} \ENG{workflows}. \\
\hline
\end{tabular}
\end{chaptertablebox}
\end{table}


Έρευνες έδειξαν ότι το μέσο \ENG{LOM} \ENG{record} συμπλήρωνε μόνο 10-15 από τα 76 πεδία, κυρίως τα βασικά (\ENG{Title}, \ENG{Description}, \ENG{Keywords}, \ENG{Technical} \ENG{Format})~\cite{barker2005metadata_practice}.

\subsection{Εξέλιξη προς Απλούστερα \ENG{Schemas}}
\label{subsec:metadata_evolution}

Η εμπειρία από την εφαρμογή του \ENG{IEEE LOM} ανέδειξε στην πράξη ότι, παρότι το πρότυπο παρείχε ένα πλούσιο και εκτενές πλαίσιο περιγραφής, η αξιοποίησή του σε πραγματικά περιβάλλοντα απαιτούσε πρόσθετες προσαρμογές και απλοποιήσεις. Στο πλαίσιο αυτό, αναγνωρίστηκε πρώτα η ανάγκη για \ENG{profiling}, δηλαδή για τη δημιουργία απλούστερων υποσυνόλων του \ENG{LOM} που να ανταποκρίνονται στις απαιτήσεις συγκεκριμένων κοινοτήτων και περιπτώσεων χρήσης, χωρίς να επιβαρύνουν τη διαδικασία καταγραφής με περιττή πολυπλοκότητα. Παράλληλα, έγινε σαφές ότι σε πολλές εφαρμογές ήταν αναγκαία η υιοθέτηση \ENG{application profiles}, τα οποία συνδυάζουν στοιχεία από πολλαπλά \ENG{metadata schemas} (όπως \ENG{Dublin Core}, \ENG{LOM} και \ENG{domain-specific} επεκτάσεις), ώστε να καλύπτονται τόσο γενικές όσο και εξειδικευμένες ανάγκες περιγραφής.

Επιπλέον, η δυσκολία της χειροκίνητης συμπλήρωσης εκτενών μεταδεδομένων οδήγησε στην κατεύθυνση της \ENG{automated metadata generation}, δηλαδή στη χρήση τεχνικών μηχανικής μάθησης και αυτοματοποιημένων μεθόδων για την εξαγωγή \ENG{metadata} απευθείας από το ίδιο το περιεχόμενο. Τέλος, αναδείχθηκε και η αξία του \ENG{social tagging}, όπου \ENG{user-generated tags} και \ENG{folksonomies} μπορούν να λειτουργήσουν συμπληρωματικά προς τα \ENG{formalized metadata}, προσφέροντας πιο ευέλικτες και «ζωντανές» περιγραφές που αντανακλούν τον τρόπο με τον οποίο οι χρήστες αντιλαμβάνονται και αναζητούν το περιεχόμενο.

\subsection{Σύγχρονη Κατάσταση \ENG{Metadata} \ENG{Standards} (2026)}
\label{subsec:metadata_modern}

Σήμερα, η κατάσταση των \ENG{metadata standards} χαρακτηρίζεται από μια πιο πολυδιάστατη και «μεικτή» πραγματικότητα, όπου δεν υπάρχει ένα μοναδικό καθολικό σχήμα που να καλύπτει όλες τις ανάγκες. Αντίθετα, παρατηρείται συνύπαρξη πολλαπλών \ENG{schemas}, τα οποία χρησιμοποιούνται ανάλογα με το πλαίσιο: το \ENG{Dublin Core} συναντάται συχνά σε πιο απλά \ENG{use cases}, τα \ENG{LOM profiles} αξιοποιούνται κυρίως σε εκπαιδευτικά \ENG{repositories}, ενώ το \ENG{Schema.org} έχει κερδίσει έδαφος λόγω της συμβολής του στη \ENG{web discoverability}. Παράλληλα, το \ENG{LRMI} (\ENG{Learning Resource Metadata Initiative}) λειτουργεί ως γέφυρα ενσωμάτωσης εκπαιδευτικών περιγραφών μέσα στο οικοσύστημα του \ENG{Schema.org}, διευκολύνοντας την προβολή και την αναζήτηση εκπαιδευτικών πόρων στο διαδίκτυο~\cite{lrmi2024guide}. Σε πιο λειτουργικό επίπεδο, τα \ENG{xAPI} \ENG{profiles} χρησιμοποιούνται για πιο λεπτομερή καταγραφή εμπειριών μάθησης, συμπληρώνοντας τα παραδοσιακά \ENG{metadata}.

Ταυτόχρονα, τα \ENG{metadata} συνδέονται ολοένα και περισσότερο με τεχνικές \ENG{AI}, καθώς αυξάνεται η ανάγκη για αυτοματοποίηση σε περιβάλλοντα με μεγάλο όγκο περιεχομένου. Σε αυτό το πλαίσιο, εφαρμόζεται \ENG{automated content analysis} για \ENG{metadata extraction}, ενώ τεχνικές \ENG{Natural Language Processing} χρησιμοποιούνται για ταξινόμηση και κατηγοριοποίηση. Επιπλέον, εργαλεία \ENG{Computer Vision} αξιοποιούνται για την περιγραφή εικόνων και \ENG{videos}, προσθέτοντας νέες δυνατότητες εμπλουτισμού μεταδεδομένων που παλαιότερα απαιτούσαν αποκλειστικά χειροκίνητη εργασία.

Μια ακόμη σύγχρονη τάση είναι η προσέγγιση του \ENG{Linked Data}, όπου η αναπαράσταση μεταδεδομένων με \ENG{RDF} και η αξιοποίηση \ENG{SPARQL} επιτρέπουν πιο δυναμική διασύνδεση και αναζήτηση πληροφορίας. Μέσω της σύνδεσης με \ENG{external vocabularies} (όπως \ENG{DBpedia} και \ENG{Wikidata}), επιδιώκεται η ενίσχυση της \ENG{semantic interoperability}, ώστε διαφορετικά συστήματα να μπορούν να «καταλαβαίνουν» με μεγαλύτερη συνέπεια το ίδιο περιεχόμενο και τις ιδιότητές του.

Παρά τις προκλήσεις, η θεμελίωση των \ENG{metadata standards} κατά την περίοδο 1995–2002 αποδείχθηκε κρίσιμη για τη διαμόρφωση και την εξέλιξη του \ENG{e-learning ecosystem}. Οι βασικές έννοιες που εισήχθησαν τότε, όπως οι \ENG{learning resource types}, το \ENG{educational context} και τα \ENG{technical requirements}, παραμένουν επίκαιρες μέχρι σήμερα, ακόμη κι αν πλέον εφαρμόζονται σε πιο απλοποιημένες ή προσαρμοσμένες μορφές.

\begin{figure}[H]
  \centering
  \resizebox{\textwidth}{!}{%
    % FIGURE 2.4: IEEE LOM (Learning Object Metadata) Structure
% Hierarchical tree showing 9 categories
\begin{chapterfigurebox}
\begin{tikzpicture}[
  scale=0.85,
  every node/.style={font=\small},
  level 1/.style={sibling distance=3.5cm, level distance=3cm},
  level 2/.style={sibling distance=2cm, level distance=2.5cm},
  root/.style={rectangle, rounded corners, draw, very thick, minimum width=4cm, minimum height=1.2cm, align=center, fill=metadatagray!30, font=\bfseries\large},
  category/.style={rectangle, rounded corners, draw, thick, minimum width=2.8cm, minimum height=1cm, align=center, fill=white, font=\small\bfseries},
  element/.style={rectangle, draw, minimum width=2.2cm, minimum height=0.6cm, align=center, font=\footnotesize}
]

% Root node
\node[root] at (0, 8) (lom) {\ENG{IEEE} \ENG{LOM}\\(\ENG{Learning} \ENG{Object} \ENG{Metadata})};
\node[above, font=\scriptsize] at (lom.north) {\ENG{IEEE} 1484.12.1-2002};

% Category 1: General
\node[category, fill=blue!20] at (-8, 5) (general) {1. \ENG{General}};
\node[element, fill=blue!10, below=0.3cm of general] {\ENG{Title}, \ENG{Language}, \ENG{Description}};

% Category 2: Lifecycle
\node[category, fill=green!20] at (-5, 5) (lifecycle) {2. \ENG{Lifecycle}};
\node[element, fill=green!10, below=0.3cm of lifecycle, text width=2.5cm] {\ENG{Version}, \ENG{Status}, \ENG{Contributors}};

% Category 3: Meta-Metadata
\node[category, fill=purple!20, text width=2.5cm] at (-2, 5) (metameta) {3. \ENG{Meta-}\\\ENG{Metadata}};
\node[element, fill=purple!10, below=0.3cm of metameta, text width=2.5cm] {\ENG{Catalog}, \ENG{Entry}, \ENG{Schema}};

% Category 4: Technical
\node[category, fill=orange!20] at (2, 5) (technical) {4. \ENG{Technical}};
\node[element, fill=orange!10, below=0.3cm of technical, text width=2.5cm] {\ENG{Format}, \ENG{Size}, \ENG{Location}};

% Category 5: Educational
\node[category, fill=red!20] at (5, 5) (educational) {5. \ENG{Educational}};
\node[element, fill=red!10, below=0.3cm of educational, text width=2.5cm] {\ENG{LearningResourceType}, \ENG{InteractivityLevel}};

% Category 6: Rights
\node[category, fill=brown!20] at (8, 5) (rights) {6. \ENG{Rights}};
\node[element, fill=brown!10, below=0.3cm of rights, text width=2.5cm] {\ENG{Cost}, \ENG{Copyright}, \ENG{Terms}};

% Category 7: Relation
\node[category, fill=cyan!20] at (-6, 1.5) (relation) {7. \ENG{Relation}};
\node[element, fill=cyan!10, below=0.3cm of relation, text width=2.5cm] {\ENG{Kind}, \ENG{Resource} (\ENG{other} \ENG{LOs})};

% Category 8: Annotation
\node[category, fill=yellow!40] at (-1, 1.5) (annotation) {8. \ENG{Annotation}};
\node[element, fill=yellow!20, below=0.3cm of annotation, text width=2.5cm] {\ENG{Person}, \ENG{Date}, \ENG{Description}};

% Category 9: Classification
\node[category, fill=violet!20] at (4, 1.5) (classification) {9. \ENG{Classification}};
\node[element, fill=violet!10, below=0.3cm of classification, text width=2.5cm] {\ENG{Purpose}, \ENG{TaxonPath}, \ENG{Keyword}};

% Connections from root to categories
\draw[thick, metadatagray] (lom) -- (general);
\draw[thick, metadatagray] (lom) -- (lifecycle);
\draw[thick, metadatagray] (lom) -- (metameta);
\draw[thick, metadatagray] (lom) -- (technical);
\draw[thick, metadatagray] (lom) -- (educational);
\draw[thick, metadatagray] (lom) -- (rights);
\draw[thick, metadatagray] (lom) -- (relation);
\draw[thick, metadatagray] (lom) -- (annotation);
\draw[thick, metadatagray] (lom) -- (classification);

% XML example box
\node[rectangle, draw, thick, fill=gray!10, align=left, font=\ttfamily\tiny, text width=6cm] at (-9, -1.5) (xml) {
\textbf{\ENG{XML} \ENG{Example} (\ENG{snippet}):}\\
<\ENG{lom} \ENG{xmlns}="...">\\
\quad <\ENG{general}>\\
\qquad <\ENG{title}>\ENG{Course} 101</\ENG{title}>\\
\qquad <\ENG{language}>\ENG{en}</\ENG{language}>\\
\quad </\ENG{general}>\\
\quad <\ENG{technical}>\\
\qquad <\ENG{format}>\ENG{text/html}</\ENG{format}>\\
\qquad <\ENG{size}>1048576</\ENG{size}>\\
\quad </\ENG{technical}>\\
</\ENG{lom}>
};

% Key statistics box
\node[rectangle, draw, thick, fill=solutiongreen!10, align=left, font=\scriptsize, text width=5cm] at (4, -1.5) {
  \textbf{\ENG{IEEE} \ENG{LOM} \ENG{Statistics}:}\\
  \textbullet~\textbf{76 \ENG{data} \ENG{elements}} \ENG{total}\\
  \textbullet~\textbf{9 \ENG{categories}} (\ENG{hierarchical})\\
  \textbullet~\ENG{XML} \ENG{Schema-based}\\
  \textbullet~\ENG{Multilingual} \ENG{support}\\
  \textbullet~\ENG{Extensible} \ENG{vocabulary}
};

% Legend
\node[below, font=\tiny, align=center] at (0, -3.5) {
  \textbf{Σημείωση:} \ENG{IEEE} \ENG{LOM} χρησιμοποιείται από \ENG{SCORM} 1.2/2004, \ENG{IMS} \ENG{CC}, και άλλα \ENG{standards}\\
  \textbf{Εναλλακτικά:} \ENG{Dublin} \ENG{Core} (15 \ENG{elements}, \ENG{simpler}), \ENG{Schema.org} (\ENG{web-focused})
};

% Adoption note
\node[below, font=\scriptsize, align=center, text width=12cm] at (0, -4.5) {
  \textbf{\ENG{Adoption}:} Υιοθετήθηκε από \ENG{ADL} (\ENG{SCORM}), \ENG{IMS}, \ENG{European} \ENG{Schoolnet}, \ENG{IEEE} \ENG{LTSC}\\
  \textbf{Κληρονομιά:} Αν και πολύπλοκο, το \ENG{LOM} θεμελίωσε βασικές έννοιες \ENG{metadata} για \ENG{learning} \ENG{objects}
};

\end{tikzpicture}
\end{chapterfigurebox}

  }
  \caption{Δομή \ENG{IEEE} \ENG{LOM} - 9 Κατηγορίες Μεταδεδομένων}
  \label{fig:ieee_lom}
\end{figure}
